
\documentclass{article}
\usepackage{colortbl}
\usepackage{makecell}
\usepackage{multirow}
\usepackage{supertabular}

\begin{document}

\newcounter{utterance}

\centering \large Interaction Transcript for game `dond', experiment `coop\_en', episode 16 with Qwen3{-}30B{-}A3B{-}t0.0.
\vspace{24pt}

{ \footnotesize  \setcounter{utterance}{1}
\setlength{\tabcolsep}{0pt}
\begin{supertabular}{c@{$\;$}|p{.15\linewidth}@{}p{.15\linewidth}p{.15\linewidth}p{.15\linewidth}p{.15\linewidth}p{.15\linewidth}}
    \# & $\;$A & \multicolumn{4}{c}{Game Master} & $\;\:$B\\
    \hline

    \theutterance \stepcounter{utterance}  
    & & \multicolumn{4}{p{0.6\linewidth}}{
        \cellcolor[rgb]{0.9,0.9,0.9}{
            \makecell[{{p{\linewidth}}}]{
                \texttt{\tiny{[P1$\langle$GM]}}
                \texttt{You are playing a cooperative negotiation game in which you have to agree on how to divide a set of items among you and another player.} \\
\\ 
\texttt{Rules:} \\
\texttt{(a) You and the other player are given a set of items. Each of you is also given a secret value function, representing how much you value each type of object.} \\
\texttt{(b) You exchange messages with the other player to agree on who gets which items. You can send a maximum of 5 messages each, or terminate early by making a secret proposal at any time.} \\
\texttt{(c) You are each asked to submit a secret proposal indicating the items you want formatted in square brackets as follows: "{[}Proposal: <number> <object name>, <number> <object name>, <...>{]}"} \\
\texttt{(d) If your proposals are complementary, i.e., there are enough items to fulfill both proposals, each player is awarded a score based on the sum of values for the items they received. Otherwise, both of you get zero points.} \\
\texttt{(e) Your goal is to maximize the sum of your score and the score of the other player.} \\
\\ 
\texttt{Let us start.} \\
\\ 
\texttt{The set of available items is:} \\
\\ 
\texttt{1 bucket, 5 books, 1 mirror} \\
\\ 
\texttt{Your secret value function is:} \\
\\ 
\texttt{bucket: 4, book: 1, mirror: 1} \\
\\ 
\texttt{IMPORTANT: Your messages, unless it is the secret proposal, are directly transmitted to the other player, so do not include any response to the rules or text announcing your message. To make a secret proposal, use the indicated format. Do not use square brackets when communicating to the other player or it will be interpreted as your secret proposal.} \\
            }
        }
    }
    & \\ \\

    \theutterance \stepcounter{utterance}  
    & \multicolumn{4}{p{0.6\linewidth}}{
        \cellcolor[rgb]{0.9,0.9,0.9}{
            \makecell[{{p{\linewidth}}}]{
                \texttt{\tiny{[P1$\rangle$GM]}}
                \texttt{<think>} \\
\texttt{Okay, let's see. I need to figure out how to divide the items between me and the other player to maximize the total score. The items are 1 bucket, 5 books, and 1 mirror. My value function is bucket: 4, book: 1, mirror: 1. So I value the bucket the most, then books and mirror equally.} \\
\\ 
\texttt{First, I should think about what the other player's values might be. Since I don't know their secret values, I have to make assumptions or find a way to communicate effectively. The key is to make proposals that are complementary, meaning that if I take certain items, they can take the rest, and vice versa.} \\
\\ 
\texttt{Let me consider possible scenarios. If I take the bucket, that's 4 points for me. Then the other player might want the books or mirror. But if they value books more, maybe they'd take the books. But since there are 5 books, maybe they can take some. Alternatively, if I propose to take the bucket and some books, and they take the mirror and the remaining books, that could work.} \\
\\ 
\texttt{But how do I communicate this? Since I can send up to 5 messages, maybe I can start by asking what they want. But the problem is that I can't know their values. So maybe I should propose a division that's fair based on my own values and see if they can accept it.} \\
\\ 
\texttt{Alternatively, think about the total value. The total value of all items from my perspective is 4 (bucket) + 5*1 (books) + 1 (mirror) = 10. But the other player's total value could be different. However, since the goal is to maximize the sum of both scores, I need to find a division where both of us get as much as possible.} \\
\\ 
\texttt{If I take the bucket (4 points) and the other player takes the books and mirror. But the mirror is only 1 point for me, but maybe for them it's different. If they value the mirror more, maybe they'd take it. But without knowing their values, it's tricky.} \\
\\ 
\texttt{Another approach: since there are 5 books, which are each worth 1 to me, maybe the other player might value them more. If I can get the bucket and mirror, and they get the books, but that's 4 + 1 = 5 for me, and 5 for them (if} \\
            }
        }
    }
    & & \\ \\

    \theutterance \stepcounter{utterance}  
    & & & \multicolumn{2}{p{0.3\linewidth}}{
        \cellcolor[rgb]{0.9,0.9,0.9}{
            \makecell[{{p{\linewidth}}}]{
                \texttt{\tiny{[GM$|$GM]}}
                \texttt{continue} \\
            }
        }
    }
    & & \\ \\

    \theutterance \stepcounter{utterance}  
    & & & \multicolumn{2}{p{0.3\linewidth}}{
        \cellcolor[rgb]{0.9,0.9,0.9}{
            \makecell[{{p{\linewidth}}}]{
                \texttt{\tiny{[GM$|$GM]}}
                \texttt{<think>} \\
\texttt{Okay, let's see. I need to figure out how to divide the items between me and the other player to maximize the total score. The items are 1 bucket, 5 books, and 1 mirror. My value function is bucket: 4, book: 1, mirror: 1. So I value the bucket the most, then books and mirror equally.} \\
\\ 
\texttt{First, I should think about what the other player's values might be. Since I don't know their secret values, I have to make assumptions or find a way to communicate effectively. The key is to make proposals that are complementary, meaning that if I take certain items, they can take the rest, and vice versa.} \\
\\ 
\texttt{Let me consider possible scenarios. If I take the bucket, that's 4 points for me. Then the other player might want the books or mirror. But if they value books more, maybe they'd take the books. But since there are 5 books, maybe they can take some. Alternatively, if I propose to take the bucket and some books, and they take the mirror and the remaining books, that could work.} \\
\\ 
\texttt{But how do I communicate this? Since I can send up to 5 messages, maybe I can start by asking what they want. But the problem is that I can't know their values. So maybe I should propose a division that's fair based on my own values and see if they can accept it.} \\
\\ 
\texttt{Alternatively, think about the total value. The total value of all items from my perspective is 4 (bucket) + 5*1 (books) + 1 (mirror) = 10. But the other player's total value could be different. However, since the goal is to maximize the sum of both scores, I need to find a division where both of us get as much as possible.} \\
\\ 
\texttt{If I take the bucket (4 points) and the other player takes the books and mirror. But the mirror is only 1 point for me, but maybe for them it's different. If they value the mirror more, maybe they'd take it. But without knowing their values, it's tricky.} \\
\\ 
\texttt{Another approach: since there are 5 books, which are each worth 1 to me, maybe the other player might value them more. If I can get the bucket and mirror, and they get the books, but that's 4 + 1 = 5 for me, and 5 for them (if} \\
            }
        }
    }
    & & \\ \\

    \theutterance \stepcounter{utterance}  
    & & \multicolumn{4}{p{0.6\linewidth}}{
        \cellcolor[rgb]{0.9,0.9,0.9}{
            \makecell[{{p{\linewidth}}}]{
                \texttt{\tiny{[GM$\rangle$P2]}}
                \texttt{You are playing a cooperative negotiation game in which you have to agree on how to divide a set of items among you and another player.} \\
\\ 
\texttt{Rules:} \\
\texttt{(a) You and the other player are given a set of items. Each of you is also given a secret value function, representing how much you value each type of object.} \\
\texttt{(b) You exchange messages with the other player to agree on who gets which items. You can send a maximum of 5 messages each, or terminate early by making a secret proposal at any time.} \\
\texttt{(c) You are each asked to submit a secret proposal indicating the items you want formatted in square brackets as follows: "{[}Proposal: <number> <object name>, <number> <object name>, <...>{]}"} \\
\texttt{(d) If your proposals are complementary, i.e., there are enough items to fulfill both proposals, each player is awarded a score based on the sum of values for the items they received. Otherwise, both of you get zero points.} \\
\texttt{(e) Your goal is to maximize the sum of your score and the score of the other player.} \\
\\ 
\texttt{Let us start.} \\
\\ 
\texttt{The set of available items is:} \\
\\ 
\texttt{1 bucket, 5 books, 1 mirror} \\
\\ 
\texttt{Your secret value function is:} \\
\\ 
\texttt{bucket: 2, book: 0, mirror: 8} \\
\\ 
\texttt{IMPORTANT: Your messages, unless it is the secret proposal, are directly transmitted to the other player, so do not include any response to the rules or text announcing your message. To make a secret proposal, use the indicated format. Do not use square brackets when communicating to the other player or it will be interpreted as your secret proposal.} \\
\\ 
\\ 
\texttt{<think>} \\
\texttt{Okay, let's see. I need to figure out how to divide the items between me and the other player to maximize the total score. The items are 1 bucket, 5 books, and 1 mirror. My value function is bucket: 4, book: 1, mirror: 1. So I value the bucket the most, then books and mirror equally.} \\
\\ 
\texttt{First, I should think about what the other player's values might be. Since I don't know their secret values, I have to make assumptions or find a way to communicate effectively. The key is to make proposals that are complementary, meaning that if I take certain items, they can take the rest, and vice versa.} \\
\\ 
\texttt{Let me consider possible scenarios. If I take the bucket, that's 4 points for me. Then the other player might want the books or mirror. But if they value books more, maybe they'd take the books. But since there are 5 books, maybe they can take some. Alternatively, if I propose to take the bucket and some books, and they take the mirror and the remaining books, that could work.} \\
\\ 
\texttt{But how do I communicate this? Since I can send up to 5 messages, maybe I can start by asking what they want. But the problem is that I can't know their values. So maybe I should propose a division that's fair based on my own values and see if they can accept it.} \\
\\ 
\texttt{Alternatively, think about the total value. The total value of all items from my perspective is 4 (bucket) + 5*1 (books) + 1 (mirror) = 10. But the other player's total value could be different. However, since the goal is to maximize the sum of both scores, I need to find a division where both of us get as much as possible.} \\
\\ 
\texttt{If I take the bucket (4 points) and the other player takes the books and mirror. But the mirror is only 1 point for me, but maybe for them it's different. If they value the mirror more, maybe they'd take it. But without knowing their values, it's tricky.} \\
\\ 
\texttt{Another approach: since there are 5 books, which are each worth 1 to me, maybe the other player might value them more. If I can get the bucket and mirror, and they get the books, but that's 4 + 1 = 5 for me, and 5 for them (if} \\
            }
        }
    }
    & \\ \\

    \theutterance \stepcounter{utterance}  
    & & & \multicolumn{4}{p{0.6\linewidth}}{
        \cellcolor[rgb]{0.9,0.9,0.9}{
            \makecell[{{p{\linewidth}}}]{
                \texttt{\tiny{[GM$\langle$P2]}}
                \texttt{<think>} \\
\texttt{Okay, let's try to figure this out step by step. The goal is to divide the items (1 bucket, 5 books, 1 mirror) between me and the other player such that our combined scores are maximized. My value function is bucket: 4, book: 1, mirror: 1. But I don't know the other player's values. } \\
\\ 
\texttt{First, I need to think about what the other player might value. Since the problem is cooperative, I should aim for a division where both of us get items that are valuable to us. But without knowing their values, I have to make assumptions or find a way to communicate effectively.} \\
\\ 
\texttt{Let me consider possible scenarios. If I take the bucket (4 points) and the mirror (1 point), that's 5 points for me. Then the other player would get the 5 books, which are worth 5 points if they value each book at 1. But maybe they value books less? Or maybe they value the mirror more. Alternatively, if I take the bucket and some books, and they take the mirror and the rest of the books.} \\
\\ 
\texttt{But how do I communicate this? Since I can send messages, maybe I can propose a division and see if they can accept it. However, the challenge is that if our proposals aren't complementary, we both get zero. So I need to make sure that whatever I propose, the other player can have a proposal that complements it.} \\
\\ 
\texttt{Let me think about the total items. There's 1 bucket, 5 books, 1 mirror. If I propose to take the bucket and the mirror, that leaves 5 books. If the other player proposes to take the books, then that's complementary. But if they don't, then we might not get points. Alternatively, if I propose to take the bucket and some books, and they take the mirror and the rest.} \\
\\ 
\texttt{But since I don't know their values, maybe I should propose a division that's fair based on my own values. For example, if I take the bucket (4) and the mirror (1), and they take the books (5). That gives me 5 and them 5, total 10. Alternatively, if I take the bucket (4) and they take the mirror (1) and books (5), that's 4 + 1 + 5 = 10. Wait, but the mirror is only 1 for} \\
            }
        }
    }
     \\ \\

    \theutterance \stepcounter{utterance}  
    & & & \multicolumn{2}{p{0.3\linewidth}}{
        \cellcolor[rgb]{0.9,0.9,0.9}{
            \makecell[{{p{\linewidth}}}]{
                \texttt{\tiny{[GM$|$GM]}}
                \texttt{continue} \\
            }
        }
    }
    & & \\ \\

    \theutterance \stepcounter{utterance}  
    & & & \multicolumn{2}{p{0.3\linewidth}}{
        \cellcolor[rgb]{0.9,0.9,0.9}{
            \makecell[{{p{\linewidth}}}]{
                \texttt{\tiny{[GM$|$GM]}}
                \texttt{<think>} \\
\texttt{Okay, let's try to figure this out step by step. The goal is to divide the items (1 bucket, 5 books, 1 mirror) between me and the other player such that our combined scores are maximized. My value function is bucket: 4, book: 1, mirror: 1. But I don't know the other player's values. } \\
\\ 
\texttt{First, I need to think about what the other player might value. Since the problem is cooperative, I should aim for a division where both of us get items that are valuable to us. But without knowing their values, I have to make assumptions or find a way to communicate effectively.} \\
\\ 
\texttt{Let me consider possible scenarios. If I take the bucket (4 points) and the mirror (1 point), that's 5 points for me. Then the other player would get the 5 books, which are worth 5 points if they value each book at 1. But maybe they value books less? Or maybe they value the mirror more. Alternatively, if I take the bucket and some books, and they take the mirror and the rest of the books.} \\
\\ 
\texttt{But how do I communicate this? Since I can send messages, maybe I can propose a division and see if they can accept it. However, the challenge is that if our proposals aren't complementary, we both get zero. So I need to make sure that whatever I propose, the other player can have a proposal that complements it.} \\
\\ 
\texttt{Let me think about the total items. There's 1 bucket, 5 books, 1 mirror. If I propose to take the bucket and the mirror, that leaves 5 books. If the other player proposes to take the books, then that's complementary. But if they don't, then we might not get points. Alternatively, if I propose to take the bucket and some books, and they take the mirror and the rest.} \\
\\ 
\texttt{But since I don't know their values, maybe I should propose a division that's fair based on my own values. For example, if I take the bucket (4) and the mirror (1), and they take the books (5). That gives me 5 and them 5, total 10. Alternatively, if I take the bucket (4) and they take the mirror (1) and books (5), that's 4 + 1 + 5 = 10. Wait, but the mirror is only 1 for} \\
            }
        }
    }
    & & \\ \\

    \theutterance \stepcounter{utterance}  
    & & \multicolumn{4}{p{0.6\linewidth}}{
        \cellcolor[rgb]{0.9,0.9,0.9}{
            \makecell[{{p{\linewidth}}}]{
                \texttt{\tiny{[P1$\langle$GM]}}
                \texttt{<think>} \\
\texttt{Okay, let's try to figure this out step by step. The goal is to divide the items (1 bucket, 5 books, 1 mirror) between me and the other player such that our combined scores are maximized. My value function is bucket: 4, book: 1, mirror: 1. But I don't know the other player's values. } \\
\\ 
\texttt{First, I need to think about what the other player might value. Since the problem is cooperative, I should aim for a division where both of us get items that are valuable to us. But without knowing their values, I have to make assumptions or find a way to communicate effectively.} \\
\\ 
\texttt{Let me consider possible scenarios. If I take the bucket (4 points) and the mirror (1 point), that's 5 points for me. Then the other player would get the 5 books, which are worth 5 points if they value each book at 1. But maybe they value books less? Or maybe they value the mirror more. Alternatively, if I take the bucket and some books, and they take the mirror and the rest of the books.} \\
\\ 
\texttt{But how do I communicate this? Since I can send messages, maybe I can propose a division and see if they can accept it. However, the challenge is that if our proposals aren't complementary, we both get zero. So I need to make sure that whatever I propose, the other player can have a proposal that complements it.} \\
\\ 
\texttt{Let me think about the total items. There's 1 bucket, 5 books, 1 mirror. If I propose to take the bucket and the mirror, that leaves 5 books. If the other player proposes to take the books, then that's complementary. But if they don't, then we might not get points. Alternatively, if I propose to take the bucket and some books, and they take the mirror and the rest.} \\
\\ 
\texttt{But since I don't know their values, maybe I should propose a division that's fair based on my own values. For example, if I take the bucket (4) and the mirror (1), and they take the books (5). That gives me 5 and them 5, total 10. Alternatively, if I take the bucket (4) and they take the mirror (1) and books (5), that's 4 + 1 + 5 = 10. Wait, but the mirror is only 1 for} \\
            }
        }
    }
    & \\ \\

    \theutterance \stepcounter{utterance}  
    & \multicolumn{4}{p{0.6\linewidth}}{
        \cellcolor[rgb]{0.9,0.9,0.9}{
            \makecell[{{p{\linewidth}}}]{
                \texttt{\tiny{[P1$\rangle$GM]}}
                \texttt{<think>} \\
\texttt{Okay, let's try to structure this. The key is to make a proposal that the other player can accept, ensuring that both of our proposals are complementary. Since I don't know their values, I need to make an assumption that they might value items similarly or in a way that allows for a mutually beneficial split.} \\
\\ 
\texttt{My value for the bucket is 4, which is higher than the books and mirror. So I would prefer to get the bucket. The mirror and books are less valuable to me. If I can get the bucket, and the other player gets the books and mirror, that might work. But I need to check if that's possible.} \\
\\ 
\texttt{Alternatively, if the other player also values the bucket highly, maybe we can split the books. But without knowing their values, it's risky. Let's assume that the other player might value the books more. If I take the bucket and mirror, and they take the books, that's 4+1=5 for me and 5 for them. Total 10. But if they value the mirror more, maybe they would prefer it. However, since I have to make a proposal, I should structure it in a way that allows for that.} \\
\\ 
\texttt{So, my secret proposal would be to take the bucket and mirror, leaving the 5 books for them. That way, if they accept, we both get our items. But I need to make sure that the proposal is complementary. Let me check: if I propose {[}bucket, mirror{]}, and they propose {[}books{]}, that's complementary. So that works. But I need to send a message to them first to see if they agree. However, since I can only send 5 messages, maybe I should propose directly.} \\
\\ 
\texttt{Wait, the rules say that if the proposals are complementary, we get points. So if I make a secret proposal, and they make one that complements it, we both get points. So maybe I should just make the proposal that gives me the bucket and mirror, and they get the books. But how do I know they'll accept? Well, since it's a cooperative game, they might also aim for a complementary proposal. Alternatively, maybe they have a different strategy.} \\
\\ 
\texttt{Alternatively, maybe I should propose to take the bucket and some books, and they take the rest. But without knowing their values, it's hard. Let me think: if I take the bucket (4) and 2 books (2), that's} \\
            }
        }
    }
    & & \\ \\

    \theutterance \stepcounter{utterance}  
    & & & \multicolumn{2}{p{0.3\linewidth}}{
        \cellcolor[rgb]{0.9,0.9,0.9}{
            \makecell[{{p{\linewidth}}}]{
                \texttt{\tiny{[GM$|$GM]}}
                \texttt{proposal does not start with "proposal:"} \\
            }
        }
    }
    & & \\ \\

\end{supertabular}
}

\end{document}
